Between late 2019 and early 2020, the red supergiant Betelgeuse experienced an
unusual visual fading event known as the ``Great Dimming'' (Figure T04.1). One
explanation is that the star ejected material that formed a dust clump,
obscuring the stellar disk. Assume the dimming occurred at a normal pulsation
minimum and was caused entirely by the cylindrical dust clump covering 80\% of
the stellar disk.

\ioaafig{2026_TH_T04-f1.pdf}

\begin{description}
\item[(a)] (2.5 points) Estimate the ratio ($k_1$) of the observed brightness
  at the minimum light to that at the normal pulsation minimum.
\item[(b)] (2.5 points) Estimate the fraction ($k_2$) of light transmitted
  through the obscured region.
\item[(c)] (1.5 points) Calculate the optical depth $\tau$ of the dust clump.
\item[(d)] (3.5 points) Assuming a mass absorption coefficient of the dust
  $\kappa = 1000.\,\mathrm{m^2\,kg^{-1}}$, a gas-to-dust mass ratio of 200.,
  and a radius of Betelgeuse of $1000\,\mathrm{R}_\odot$, estimate the total
  mass ($M_{\mathrm{eject}}$) of the ejected material in units of
  $\mathrm{M}_\odot$.
\end{description}
